\section{The formal verification}\label{app:formal}

The proof of this paper is formalized in Lean~4~\cite{Lean4} with Mathlib~\cite{Mathlib20}, in the
repository~\cite{Companion}. This appendix says what the formalization states, how it is checked,
on what it depends, and where its arguments differ from the text, and it lists the Lean
declarations that prove the numbered results of the paper.

\subsection{The statements}\label{app:formal-statements}

The file \texttt{Palomar/Challenge.lean} imports only Mathlib and states the three parts of
\cref{thm:main} as the theorems \lean{octagon_density}, \lean{reinhardt} and
\lean{reinhardt_equality} of the namespace \lean{ReinhardtChallenge}. Its definitions are those of
\cref{def:packing-density}, which is Definition~2.1.1 of~\cite{HalesVajjha2024}. Points of the
plane are maps $\{0,1\}\to\R$ and areas are Lebesgue measure. A convex disk is a compact convex set
with nonempty interior, and a congruent copy of $K$ is its image under a map $x\mapsto Ax+b$ with
$A$ orthogonal. A packing is a family of convex disks whose members have pairwise disjoint
interiors. Its density exists when the upper and lower limits, as $r\to\infty$, of the fraction of
the disk of radius $r$ about the origin that it covers are finite and equal, and $\delta(K)$ is the
supremum of the densities that exist over the packings by congruent copies of $K$. A disk is
centrally symmetric when there is a point $c$ with $2c-x\in K$ for every $x\in K$. The smoothed
octagon is given in Reinhardt's description (\cref{def:smoothed-octagon}): with $d_k(x)$ the signed
distance from $x$ to the line of the $k$-th side of the regular octagon of inradius~$1$, it is the
octagon $d_k\ge0$ with the corner between sides $k$ and $k+1$ cut off along the hyperbola
$d_{k-1}d_{k+2}=\sqrt2/4$. These definitions are copied word for word from the library, so the
theorems the library proves are the theorems the file states.

\subsection{The check}\label{app:formal-check}

The root module \texttt{SmoothedOctagon.lean} restates the three theorems and proves them with no
hypothesis. Lean's kernel checks every step, and every theorem depends only on Lean's standard
axioms \lean{propext}, \lean{Classical.choice} and \lean{Quot.sound}; an audit script checks this
for each of the 8{,}294 declarations of the library. The Comparator configuration
\texttt{Palomar/comparator.json} checks that the root module proves the statements of
\texttt{Palomar/Challenge.lean} with the same definitions. The proof is registered in the Palomar
registry~\cite{Palomar} as \texttt{PALOMAR-2026-11-09-000008}, at trust level high: Palomar's
verifier checked it with the independent kernels nanoda and con-ron at commit
\texttt{9ba6bf8} of~\cite{Companion}. Later commits rename and reorganize the library and leave the
three statements unchanged; the names below are those of the current library. Lean and Mathlib are
pinned at version \texttt{v4.36.0-rc2}. The library has about 110{,}000 lines in 344 modules, and a
full build takes about 3.5 CPU-hours.

The library depends on two others. The first~\cite{SmoothedPolygonLean} formalizes the parts of
the book of Hales and Vajjha~\cite{HalesVajjha2024} that the proof uses, together with the classical
results behind them: among them Fejes T\'oth's theorem (\cref{thm:fejes-toth}), the theorems of
Minkowski and Mahler on critical lattices and circumscribed hexagons, the compact domain
(\cref{thm:compact-domain}) and the lift of a global minimizer to a Pontryagin extremal
(\cref{thm:pmp-lift}). The second~\cite{OptimalControlLean} proves the Pontryagin maximum
principle~\cite{Pontryagin1962} with free terminal time, an integral cost and mixed boundary
conditions, by the argument of Agrachev and Sachkov~\cite{AgrachevSachkov2004}, and Filippov's
theorem~\cite{Filippov1962}; Brouwer's fixed-point theorem comes from~\cite{FixedPointLean}.

\subsection{The scalar estimates}\label{app:formal-scalar}

Most scalar inequalities of the paper are proved in Lean as in the text: polynomial identities by
ring normalization, polynomial inequalities by nonlinear arithmetic with explicit products as
hints, monotonicity and concavity in one variable by derivatives, and the logarithm by the tangent
and Pad\'e bounds of the text and Mathlib's bounds for $\log2$ and $\pi$. Some polynomial signs on
boxes, in class $\cls{F}$ (\cref{sec:forward}) and in the nonshort windows of class $\cls{M2}$
(\cref{sec:mixed}), are checked by expansion in the Bernstein basis of the box: scripts in
\texttt{scripts/certificates/} of~\cite{Companion} compute the coefficients, \lean{ring} checks the
identity between the polynomial and its Bernstein form, and normalization checks the signs of the
rational coefficients. Neither interval arithmetic nor floating-point computation enters the
proof.

\subsection{Where the formal proof argues differently}\label{app:formal-routes}

The statements of the paper are those that Lean proves, sometimes in a simpler notation. In a few
places the text proves an intermediate inequality by a different argument from the Lean proof,
usually because the hand argument is shorter to read and the formal one shorter to check:
\begin{itemize}
\item In class $\cls{F}$, the four quartic signs of the conditional Casimir bound
  (\cref{lem:fwd-casimir}) have hand arguments in the text; Lean proves them by nonlinear
  arithmetic with products of Bernstein basis polynomials as hints.
\item In the nonshort windows of class $\cls{M2}$ (\cref{app:nonshort-estimates}), the text proves
  the scale strip by monotonicity in $t$ and a comparison at $t_0=13/35$, the low ratio by shape
  monotonicity, and the bound $\chi>4/7$, the elimination of $q$ and the envelope by concavity
  arguments; Lean proves the corresponding signs by Bernstein expansion, and the low ratio by
  concavity in $k$ with Bernstein expansions at both ends.
\item In class $\cls{M1}$ the text proves the two-centre bound in the exact form
  $\cost>2\sqrt6-\sqrt3$ (\cref{prop:m1-two-centre}); Lean proves the weaker numerical bound that
  the argument needs.
\item In the smaller-short window (\cref{sec:smaller-short}), the text proves the outer face sign
  of gate~S3 (\cref{prop:ss-face-S3}) on a different domain from the Lean statement; both domains
  contain the cases that occur.
\end{itemize}
The full small-scale exclusion of \cref{rem:m2-low-scale} is not formalized; the proof uses only
its unit face (\cref{lem:m2-low-unit}).

\subsection{Dictionary}\label{app:formal-dictionary}

\Cref{tab:formal-dictionary} lists, for each numbered result of the paper, the Lean declarations
that prove it. Names are relative to the namespace \lean{SmoothedOctagon}; the prefixes
\lean{ReinhardtChallenge} and \lean{SmoothedPolygons} mark the statements of the challenge file and
the declarations of~\cite{SmoothedPolygonLean}. Several names in a row are the pieces of one
result.

{\small
\begin{longtable}{@{}>{\raggedright\arraybackslash}p{0.25\textwidth}>{\raggedright\arraybackslash}p{0.71\textwidth}@{}}
\caption{The Lean declarations behind the numbered results.}\label{tab:formal-dictionary}\\
\toprule
Result & Lean declarations\\
\midrule
\endfirsthead
\toprule
Result & Lean declarations\\
\midrule
\endhead
\bottomrule
\endfoot
\multicolumn{2}{@{}l}{\emph{Sections~1 and~2}}\\[2pt]
\cref{thm:control} & \lean{reinhardt_cost}, \lean{controlLowerBound} \\
\cref{thm:main}(a) & \lean{ReinhardtChallenge.octagon_density}, \lean{smoothedOctagon_spec} \\
\cref{thm:main}(b) & \lean{ReinhardtChallenge.reinhardt}, \lean{reinhardtConjecture} \\
\cref{thm:main}(c) & \lean{ReinhardtChallenge.reinhardt_equality}, \lean{reinhardt_equality_trajectory} \\
\cref{def:packing-density} & \lean{ReinhardtChallenge.greatestPackingDensity} \\
\cref{def:smoothed-octagon} & \lean{ReinhardtChallenge.smoothedOctagon} \\
\cref{def:global-minimizer} & \lean{SmoothedPolygons.IsReinhardtMinimizer} \\
\cref{thm:compact-domain} & \lean{SmoothedPolygons.theorem5_0_2} \\
\addlinespace
\multicolumn{2}{@{}l}{\emph{Section~3}}\\[2pt]
\cref{def:pmp-extremal} & \lean{SmoothedPolygons.IsPontryaginExtremal} \\
\cref{thm:pmp-lift} & \lean{SmoothedPolygons.exists_isPontryaginExtremal_of_isReinhardtMinimizer} \\
\cref{lem:vertex-hamiltonians} & \lean{IsPontryaginExtremal.continuousOn_vertexHamiltonian}, \lean{IsPontryaginExtremal.vertexHamiltonian_max_eq_zero} \\
\cref{red:singular-locus} & \lean{hamiltonian_vertices_equal_iff_reducedCostate_eq_zero}, \lean{hamiltonian_maximizers_in_proper_edge} \\
\cref{red:winding} & \lean{Reduction.clr_winding}, \lean{Reduction.clr_angular} \\
\cref{red:equilibrium} & \lean{Reduction.clr_vertex_bounds} \\
\cref{red:budget} & \lean{Reduction.clr_budget}, \lean{Reduction.clr_lam_sq_le} \\
\cref{prop:clearance} & \lean{Reduction.covectorNormSq_le_singularGapSq}, \lean{Reduction.switch_singularGapSq_pos} \\
\cref{prop:edge-bound} & \lean{Reduction.edgeBound_of_isReinhardtMinimizer}, \lean{Reduction.Edge.edge_bangBang} \\
\cref{red:local-exclusion} & \lean{Reduction.switch_exists_excluded_vertex} \\
\cref{prop:finite-switches} & \lean{Reduction.exists_isPiecewiseVertex} \\
\cref{def:vertex-representation} & \lean{VertexRepresentation} \\
\cref{lem:word-reduction} & \lean{Reduction.exists_rebase_vertexRepresentation}, \lean{Reduction.isReinhardtMinimizer_rebase} \\
\cref{def:normal-period} & \lean{NormalPeriod}, \lean{NormalPeriod.IsCompetitive}, \lean{NormalPeriod.InCompactDomain}, \lean{NormalPeriod.activeMargin} \\
\cref{thm:normality} & \lean{Reduction.lam_lt_zero_of_vertexRepresentation}, \lean{Reduction.exists_normal_lift} \\
\cref{red:abnormal-transport} & \lean{Reduction.gn_piece} \\
\cref{thm:finite-normal} & \lean{finiteNormalRepresentation}, \lean{Reduction.cost_le_octagonCost} \\
\addlinespace
\multicolumn{2}{@{}l}{\emph{Sections~4 and~5, \cref{app:reverse-estimates}}}\\[2pt]
\cref{lem:one-arc} & \lean{vertexPiece_eq}, \lean{vertexPiece_margin}, \lean{vertexPiece_cost}, \lean{NormalPeriod.piece_margin}, \lean{NormalPeriod.cost_eq_sum} \\
\cref{lem:per-factors} & \lean{CostateCycle.factor_pp}, \lean{CostateCycle.factor_mp}, \lean{CostateCycle.factor_pm}, \lean{CostateCycle.centre_factor} \\
\cref{lem:per-signed-action} & \lean{CostateCycle.rawCost_eq}, \lean{CostateCycle.arc_eq}, \lean{NormalPeriod.cost_eq_rawCost} \\
\cref{lem:tangent-closure} & \lean{NormalPeriod.tangentProduct_N}, \lean{Term.triC_conj_tangentFactor}, \lean{Term.triC_conj_R} \\
\cref{def:signs} & \lean{VertexRepresentation.IsForward}, \lean{VertexRepresentation.IsReverse}, \lean{VertexRepresentation.IsReverseRun}, \lean{VertexRepresentation.IsForwardRun} \\
\cref{lem:sign-sum} & \lean{NormalPeriod.ff_N_mod_three} \\
\cref{def:classes} & \lean{VertexRepresentation.AllReverse}, \lean{VertexRepresentation.AllForward}, \lean{VertexRepresentation.ClassM3}, \lean{VertexRepresentation.ClassM1}, \lean{VertexRepresentation.ClassM2} \\
\cref{prop:partition} & \lean{VertexRepresentation.partition} \\
\cref{def:costate-coordinates} & \lean{costateL}, \lean{costateA}, \lean{costateB}, \lean{costateAlpha}, \lean{costateBeta}, \lean{CostateCycle.sig} \\
\cref{lem:costate-transport} & \lean{margin_mul_hamiltonian}, \lean{costateB_perp}, \lean{vertexPiece_costate}, \lean{NormalPeriod.cut_transport}, \lean{NormalPeriod.sig_le_of_mem}, \lean{NormalPeriod.cut_sig_grow_le}, \lean{NormalPeriod.cut_sig_decr_le}, \lean{CostateCycle.sig_act_succ}, \lean{NormalPeriod.costateCycle} \\
\cref{def:normal-multiplier} & \lean{NormalPeriod.Ucal}, \lean{CostateCycle.mu}, \lean{NormalPeriod.mu_zero}, \lean{CostateCycle.gam2}, \lean{NormalPeriod.det_Ucal}, \lean{triY}, \lean{CostateCycle.Mloc} \\
\cref{lem:per-switch-vectors} & \lean{CostateCycle.yv_norm}, \lean{CostateCycle.cross_pp}, \lean{CostateCycle.cross_mm}, \lean{CostateCycle.yv_pair}, \lean{CostateCycle.outv_turn}, \lean{CostateCycle.yv_add_N}, \lean{CostateCycle.pre_N} \\
\cref{lem:stationarity} & \lean{CostateCycle.stationarity}, \lean{CostateCycle.centre_stat} \\
\cref{lem:switch-norms} & \lean{CostateCycle.swd_fwd}, \lean{CostateCycle.swd_rev}, \lean{CostateCycle.turn_secant}, \lean{CostateCycle.turn_swd_ge} \\
\cref{lem:gram} & \lean{CostateCycle.gram_pp}, \lean{CostateCycle.gram_mm} \\
\cref{prop:ellipticity} & \lean{NormalPeriod.gam2_pos_of_mixed}, \lean{NormalPeriod.wire_gam2_pos_of_allForward} \\
\cref{prop:cost-casimir} & \lean{NormalPeriod.cost_casimir}, \lean{el_winding} \\
\cref{cor:casimir-cap} & \lean{NormalPeriod.gam2_le_cap} \\
\cref{thm:at-most-five} & \lean{period_le_five} \\
\cref{lem:rev-N} & \lean{cstRevN_neg} \\
\cref{lem:rev-identities} & \lean{cst_rev4_third_identity}, \lean{cst_rev4_second_identity} \\
\cref{lem:reverse-four} & \lean{CostateCycle.reverse_four_third}, \lean{CostateCycle.reverse_four_second} \\
\cref{thm:class-R} & \lean{reverse_exclusion} \\
\cref{lem:three-arc} & \lean{CostateCycle.three_arc}, \lean{CostateCycle.three_arc_not_min}, \lean{CostateCycle.turn_up}, \lean{CostateCycle.turn_down}, \lean{cst_three_grow}, \lean{cst_three_decr} \\
\cref{prop:run-structure} & \lean{reverseRun_two_and_centre_min} \\
\cref{thm:class-M3} & \lean{minimumStructure_exclusion} \\
\cref{cor:rev-mixed} & \lean{NormalPeriod.rev_run}, \lean{NormalPeriod.centre_min}, \lean{NormalPeriod.swd_ge_of_mixed} \\
\cref{lem:rev-three-gram} & \lean{el_three_gram} \\
\cref{lem:rev-cross-signs} & \lean{CostateCycle.cre_pos_pp}, \lean{CostateCycle.cre_neg_mm}, \lean{el_cross_bound} \\
\cref{lem:rev-pade} & \lean{cst_log_pade}, \lean{cst_zeta_gt}, \lean{cst_eta_lt_zeta} \\
\addlinespace
\multicolumn{2}{@{}l}{\emph{Section~6, \cref{app:short-run-estimates}}}\\[2pt]
\cref{thm:class-M1} & \lean{shortRuns_exclusion} \\
\cref{lem:m1-pattern} & \lean{VertexRepresentation.m1_pattern} \\
\cref{lem:m1-centre-data} & \lean{NormalPeriod.m1_centreData}, \lean{m1_CentreData} \\
\cref{prop:m1-angle-gap} & \lean{m1_angle_gap} \\
\cref{cor:m1-angle} & \lean{m1_arccos_gt}, \lean{m1_cos_two_pi_div_five} \\
\cref{prop:m1-winding} & \lean{m1_winding}, \lean{wind_winding_aux}, \lean{wind_polar}, \lean{wind_centre}, \lean{wind_track_range}, \lean{wind_seam} \\
\cref{lem:m1-recurrence} & \lean{CostateCycle.m1_next_centre}, \lean{CostateCycle.next_centre_local}, \lean{CostateCycle.cross_mm}, \lean{CostateCycle.pre_N} \\
\cref{lem:m1-euler} & \lean{CostateCycle.euler_sum}, \lean{CostateCycle.rawCost_eq}, \lean{elGa_euler} \\
\cref{prop:m1-two-centre} & \lean{m1_two_centre}, \lean{m1_two_centre_numeric} \\
\cref{lem:m1-log} & \lean{ag_log_le_tangent}, \lean{ag_log_le_atanh}, \lean{ag_half_sub_inv_le_log}, \lean{ag_sigma_superadd} \\
\cref{lem:m1-one-switch} & \lean{ag_phi_le} \\
\cref{lem:m1-reduce} & steps of the proof of \lean{m1_angle_gap} \\
\cref{lem:m1-core} & \lean{ag_core} \\
\cref{lem:m1-bc} & \lean{m1_bc_bound} \\
\addlinespace
\multicolumn{2}{@{}l}{\emph{Section~7, \cref{app:forward-estimates}}}\\[2pt]
\cref{thm:class-F} & \lean{forward_exclusion} \\
\cref{lem:fwd-count} & \lean{NormalPeriod.ff_N_mod_three} \\
\cref{lem:fwd-cycle} & \lean{CostateCycle.ff_act_mul_lt}, \lean{CostateCycle.ff_q_eq}, \lean{CostateCycle.ff_rawCost} \\
\cref{lem:fwd-seven} & \lean{ff_seven_le} \\
\cref{prop:fwd-window} & \lean{CostateCycle.ff_rows}, \lean{CostateCycle.ff_floors}, \lean{CostateCycle.ff_gram}, \lean{NormalPeriod.ff_gam2_eq}, \lean{ff_Win.toDim} \\
\cref{prop:fwd-spectral} & \lean{NormalPeriod.ff_cost_casimir}, \lean{NormalPeriod.wire_gam2_pos_of_allForward}, \lean{NormalPeriod.wire_spectral} \\
\cref{cor:fwd-cap} & \lean{ff_gam2_cap} \\
\cref{lem:fwd-small} & \lean{ff_Dim.U_lt_three}, \lean{fs_U_lt_three} \\
\cref{lem:fwd-threshold} & \lean{ff_Dim.eleven_fifths_lt} \\
\cref{lem:fwd-budget} & \lean{ff_Dim.budget}, \lean{ff_Dim.h_le_c}, \lean{ff_budget_g} \\
\cref{lem:fwd-neighbour} & \lean{ff_Dim.neighbor}, \lean{ff_resN_neg} \\
\cref{lem:fwd-norms} & \lean{ff_Dim.norms}, \lean{ff_Dim.norms_nonneg}, \lean{ff_resS_neg}, \lean{ff_resV_neg} \\
\cref{lem:fwd-casimir} & \lean{ff_casimir_lower} \\
\addlinespace
\multicolumn{2}{@{}l}{\emph{Section~8, \cref{app:nonshort-estimates}}}\\[2pt]
\cref{thm:class-M2} & \lean{nonshort_exclusion} \\
\cref{lem:m2-casimir} & \lean{NormalPeriod.wire_spectral}, \lean{NormalPeriod.wire_cas_lt}, \lean{NormalPeriod.dens_cas_gt} \\
\cref{lem:m2-scale} & \lean{m2_centre}, \lean{m2_scale} \\
\cref{prop:m2-small-forward} & \lean{m2_small_forward_minimum} \\
\cref{prop:m2-centre} & \lean{m2_min_centre} \\
\cref{lem:m2-reflection} & \lean{m2_reflect}, \lean{m2_reflect_fwd_iff}, \lean{m2_reflect_classM2} \\
\cref{prop:m2-window} & \lean{m2_window} \\
\cref{prop:m2-continuation} & \lean{m2_window}, \lean{M2Window.ForwardNext}, \lean{M2Window.ReverseNext} \\
\cref{lem:m2-comparison} & \lean{m2_comparison} \\
\cref{lem:m2-shape} & \lean{nsrN_mono_k}, \lean{nsrA_pos} \\
\cref{lem:m2-necessary} & \lean{nsr_window_basic}, \lean{nsr_window_N_neg}, \lean{nsr_high_t}, \lean{nsr_scale_lt}, \lean{nsr_scale_k}, \lean{nsr_centre_gap}, \lean{nsr_centre_gap_k}, \lean{nsr_window_dur}, \lean{nsr_window_dur_x}, \lean{nsr_reverse_order}, \lean{nsr_reverse_order_k} \\
\cref{prop:m2-reverse-core} & \lean{nsr_reverse_core} \\
\cref{thm:m2-reverse} & \lean{m2_reverse_exclusion} \\
\cref{prop:m2-forward-steps} & \lean{nsf_low_ratio}, \lean{nsf_small_gap_j}, \lean{nsf_small_gap} \\
\cref{thm:m2-forward} & \lean{m2_forward_exclusion} \\
\cref{lem:m2-secants} & \lean{nsr_sec_swap}, \lean{nsr_sec_start}, \lean{nsr_m2F_eq_sec}, \lean{nsr_sec_lt_chord}, \lean{nsr_turn_floor} \\
\cref{lem:m2-logs} & \lean{nsr_log_one_add_gt}, \lean{nsr_secN_lt_vbar}, \lean{nsr_log_lt_Lp}, \lean{nsr_m2F_gt_Fl}, \lean{nsf_log_gt_G}, \lean{nsf_secN_lt_VG}, \lean{nsf_secN_lt_Vq}, \lean{nsr_t_lt_T} \\
\cref{lem:m2-relax} & \lean{nsrNv_nonpos}, \lean{nsrNL_nonpos}, \lean{nsr_N_pos_of_relax} \\
\cref{lem:m2-high-t} & \lean{nsr_high_t'}, \lean{nsr_row_cubic} \\
\cref{lem:m2-unit-shape} & \lean{nsr_unit_shape}, \lean{nsr_cubic_pos}, \lean{nsrU_signs} \\
\cref{lem:m2-order-face} & \lean{nsr_order_face}, \lean{nsr_q_le_order}, \lean{nsr_concave_min} \\
\cref{lem:m2-strip} & \lean{nsr_strip} \\
\cref{lem:m2-low-unit} & \lean{nsr_low_unit}, \lean{nsr_quartic_pos} \\
\cref{lem:m2-low-ratio} & \lean{nsf_lowRatio_core}, \lean{nlr_unit_pos} \\
\cref{lem:m2-large-gap} & \lean{nsf_LGN_core}, \lean{nsf_LGB_core}, \lean{nlt_unit_pos}, \lean{nlt_shape_ge}, \lean{hrl_lower_pos}, \lean{hrs_shape_lt} \\
\cref{lem:m2-chi} & \lean{nsf_chi_core} \\
\cref{lem:m2-slope} & \lean{nsf_cas_slope} \\
\cref{lem:m2-qelim} & \lean{nsf_qelim} \\
\cref{lem:m2-envelope} & \lean{nsf_envelope} \\
\addlinespace
\multicolumn{2}{@{}l}{\emph{Section~9, \cref{app:short-estimates}}}\\[2pt]
\cref{def:ss-window} & \lean{M2Short} \\
\cref{prop:ss-extraction} & \lean{m2_short_window} \\
\eqref{eq:ss-defect} & \lean{NormalPeriod.df_window}, \lean{NormalPeriod.df_defect} \\
\cref{thm:m2-smaller-short} & \lean{m2_smallerShort_exclusion} \\
\cref{lem:ss-first} & \lean{M2Short.ss_C_le_B}, \lean{SmallerShortData.u_le_V}, \lean{M2Short.ss_u_le}, \lean{ss_order}, \lean{SmallerShortData.root_zero}, \lean{SmallerShortData.q_le_qc}, \lean{SmallerShortData.q_lt_seventh} \\
\cref{lem:ss-orders} & \lean{ss_strong_order}, \lean{ss_upper_order_54}, \lean{ss_upper_order_98}, \lean{ss_small_upper_order}, \lean{ss_quad_order}, \lean{SmallerShortData.strong_box} \\
\cref{prop:ss-S0} & \lean{M2Short.ratio_lt_twelve_fifths}, \lean{SmallerShortData.ratio_lt_twelve_fifths}, \lean{lgr_D_pos}, \lean{lgr_pairing} \\
\cref{prop:ss-S1} & \lean{M2Short.two_lt_ratio}, \lean{SmallerShortData.two_lt_ratio}, \lean{SmallerShortData.lr_cap_env} \\
\cref{prop:ss-S2} & \lean{M2Short.ratio_lt_thirteen_sixths}, \lean{SmallerShortData.ratio_lt_thirteen_sixths} \\
\cref{prop:ss-S3} & \lean{M2Short.displacement_lt}, \lean{SmallerShortData.displacement_lt} \\
\cref{prop:ss-S4} & \lean{M2Short.upper_order_casimir_floor}, \lean{SmallerShortData.upper_order_casimir_floor}, \lean{SmallerShortData.cas_cap_env} \\
\cref{prop:ss-S5} & \lean{M2Short.shear_neg}, \lean{SmallerShortData.shear_neg}, \lean{shr_residual_neg}, \lean{shr_shear_sign} \\
\cref{prop:ss-S6} & \lean{M2Short.octagonCost_lt_completion}, \lean{SmallerShortData.octagonCost_lt_completion}, \lean{SmallerShortData.cost_angle}, \lean{SmallerShortData.cost_continuation}, \lean{cost_Fc_gt}, \lean{cost_constants} \\
\cref{lem:ss-first-contact} & \lean{ss_first_contact} \\
\cref{lem:ss-outer-scale} & \lean{ss_rev_centre_scale} \\
\cref{prop:ss-face-S0} & \lean{lgr_central_tau_two}, \lean{lgr_central_duration}, \lean{lgr_outer_tests} \\
\cref{prop:ss-face-S1} & \lean{lr_central_face}, \lean{lr_outer_face} \\
\cref{prop:ss-face-S2} & \lean{ur_central_sign}, \lean{ur_outer_sign_small}, \lean{ur_outer_sign_large} \\
\cref{prop:ss-face-S3} & \lean{dsp_central_alpha_proof}, \lean{dsp_outer_face_proof} \\
\cref{prop:ss-face-S4} & \lean{cas_floor_box} \\
\cref{prop:ss-face-S6} & \lean{cost_B_concave}, \lean{cost_B_endpoint} \\
\addlinespace
\multicolumn{2}{@{}l}{\emph{Sections~10 and~11, \cref{app:terminal-estimates}}}\\[2pt]
\cref{def:term-cycle} & \lean{Term.Cycle}, \lean{Term.Cycle.rawCost}, \lean{Term.rawArc} \\
\cref{lem:term-cycle} & \lean{NormalPeriod.termCycle}, \lean{NormalPeriod.cost_eq_rawCost} \\
\cref{prop:terminal-words} & \lean{Term.Cycle.cycle_one_false}, \lean{Term.Cycle.cycle_two_false}, \lean{Term.Cycle.cycle_three_false}, \lean{Term.Cycle.cycle_four_forward}, \lean{Term.Cycle.cycle_five} \\
\cref{lem:four-closure} & \lean{Term.four_closure} \\
\cref{lem:four-param} & \lean{Term.Cycle.four_param} \\
\cref{thm:four-density} & \lean{Term.fourPiece_density_ge}, \lean{Term.fourPiece_density_eq}, \lean{ra_four_density} \\
\cref{cor:four-equality} & \lean{Term.Cycle.four_rawCost_ge}, \lean{Term.four_equality_data} \\
\cref{lem:five-closure} & \lean{Term.five_closure}, \lean{Term.Cycle.five_param} \\
\cref{lem:five-param} & \lean{Term.Cycle.five_param} \\
\cref{thm:five-density} & \lean{Term.fivePiece_density_gt} \\
\cref{thm:terminal} & \lean{terminal_cost}, \lean{Term.rawCost_ge_of_le_five} \\
\cref{prop:terminal-equality} & \lean{terminal_equality}, \lean{Term.four_of_rawCost_eq}, \lean{Term.octagon_identification} \\
\cref{thm:minimizer-cost} & \lean{minimizer_cost_ge} \\
\cref{lem:eq-realization} & \lean{eq_worstDisk_minimizer_realization} \\
\cref{lem:eq-frame-twist} & \lean{eq_frame_twist_of_contact} \\
\cref{lem:eq-frontier} & \lean{eq_disk_eq_of_frontier_eq} \\
\cref{thm:rigidity} & \lean{worstDisk_eq_linear_image} \\
\cref{prop:converse} & \lean{bridge_greatestPackingDensity_linear_image}, \lean{bridge_greatestPackingDensity_translate}, \lean{greatestPackingDensity_linear_image_octagon} \\
\cref{prop:reinhardt-construction} & \lean{cls_image_octagon}, \lean{cls_exists_affine_iff} \\
\cref{lem:log-sinh} & \lean{ra_half_sub_inv_le_log} \\
\cref{lem:four-weight} & \lean{ra_weight_monotoneOn}, \lean{ra_weight_le} \\
\cref{lem:five-h} & \lean{ra_fiveH_concaveOn}, \lean{ra_fiveH_circle_le} \\
\cref{lem:five-envelope} & \lean{ra_five_envelope}, \lean{Term.fiveEnvelope} \\
\cref{lem:five-d} & a private lemma of \lean{Term.fivePiece_density_gt} \\
\end{longtable}
}
